\documentclass[11pt]{amsart}
\usepackage{amsmath,amssymb}
\setlength{\emergencystretch}{3em}
\tolerance=600
\usepackage{tikz-cd}
\usepackage[hidelinks]{hyperref}

\newtheorem{thm}{Theorem}[section]
\newtheorem{lem}[thm]{Lemma}
\newtheorem{cor}[thm]{Corollary}
\newtheorem{prop}[thm]{Proposition}
\theoremstyle{definition}
\newtheorem{defi}[thm]{Definition}
\newtheorem{rem}[thm]{Remark}
\newtheorem{ex}[thm]{Example}
\newtheorem{que}[thm]{Question}

\newcommand{\Same}{\mathrm{Same}}
\newcommand{\Mor}{\mathrm{Mor}}
\newcommand{\Presh}{\mathrm{Presh}}
\newcommand{\Hom}{\mathrm{Hom}}
\newcommand{\im}{\mathrm{im}}
\newcommand{\Ims}{\mathrm{Im}}
\newcommand{\Eq}{\mathrm{Eq}}
\newcommand{\Sh}{\mathrm{Sh}}

\begin{document}

\title[Joins of Grothendieck topologies and the sameness lattice]{Joins of Grothendieck topologies in the lattice of two-out-of-three classes: the finite-iteration criterion}
\author{Ryota Shibusawa}
\thanks{Corresponding author. E-mail: \texttt{r-shibusawa@daiichi-koudai.ac.jp}}
\address{Daiichi Institute of Technology}
\email{r-shibusawa@daiichi-koudai.ac.jp}
\keywords{Grothendieck topology, Lawvere--Tierney topology, sheafification, local isomorphism,
  two-out-of-three, idempotent ideal, join of topologies, finite iteration}
\subjclass[2020]{18F10 (primary), 18F20, 18A32, 20M12 (secondary)}

\begin{abstract}
For a small category $S$, every Grothendieck topology $J$ defines the class
$W_J$ of maps of presheaves inverted by sheafification, an element of the
complete lattice $\Same(\Presh(S))$ of classes of morphisms closed under
composition and two-out-of-three. The map $J\mapsto W_J$ preserves all
meets; whether it preserves the join of two topologies was known to depend on
the category. We describe the join $W_J\vee W_{J'}$ for arbitrary topologies
$J,J'$: it is the class of maps whose image sieves and equaliser sieves are
reached by a bounded number of alternating applications of the closure
operators $j,j'$ of the two topologies, and every such map is a composite of
maps of $W_J\cup W_{J'}$, so the join is the closure of the union under
composition alone. Consequently $W_J\vee W_{J'}=W_{J\vee J'}$ if and only if
$(jj')^k=j\vee j'$ for some $k$, the finite-iteration axiom for Lawvere--Tierney
topologies, the same condition that governs joins of nuclei on a Heyting
algebra and on the algebra of truth values of fuzzy sets. For topologies of
ideal type, $J_D$ and $J_E$ for idempotent ideals $D,E$, the condition is the
stabilisation of the chain $DE\supseteq DEDE\supseteq\cdots$, which contains
the known positive answer for finite categories and the negative answer on
the natural numbers; for ideals generated by two sets of objects $A,B$ it
decides when the join of ``bijective at $A$'' and ``bijective at $B$'' is
``bijective at $A\cap B$''; on the natural numbers, the join of the classes of
maps bijective at the even and at the odd points is the class of maps of
towers that are bijective up to a uniform shift, a proper subclass of the
isomorphisms of pro-sets.
\end{abstract}

\maketitle

\section{Introduction}

Let $S$ be a small category and $\Presh(S)$ its category of presheaves. A
Grothendieck topology $J$ on $S$ has a sheafification functor $a_J$, and the
class $W_J$ of maps inverted by $a_J$, the $J$-local isomorphisms, is a
\emph{sameness}: a class of morphisms containing the identities and closed
under composition and two-out-of-three. Samenesses form a complete lattice
$\Same(\Presh(S))$ under inclusion, meets being intersections and the join of
two classes the least sameness containing both \cite{Same}. The assignment
$J\mapsto W_J$ preserves all meets \cite[Thm.~3.1]{Topjoin}, and the question
studied in \cite{Topjoin,Finitejoin} is whether it preserves binary joins:
whether every map inverted by the sheafification of $J\vee J'$ is obtained
from $J$-local and $J'$-local isomorphisms by finitely many compositions and
cancellations. The answer was known to be yes for every finite category
\cite{Finitejoin} and no for the poset $\mathbb N$ \cite[Thm.~7.1]{Topjoin}.
This paper describes the join $W_J\vee W_{J'}$ for arbitrary topologies and
gives the exact criterion for it to equal $W_{J\vee J'}$.

A topology $J$ has a closure operator $j$ on sieves: for a sieve $S$ on $c$,
$j(S)$ is the sieve of morphisms $b$ into $c$ with $b^*S$ covering; $S$ is
covering iff $j(S)$ is the maximal sieve. A map $f\colon X\to Y$ of presheaves
is in $W_J$ iff its \emph{image sieves} $\{b:y\cdot b\in\im f\}$, $y\in Y(c)$,
and its \emph{equaliser sieves} $\{b:x\cdot b=x'\cdot b\}$, $f(x)=f(x')$, are
covering. Say a sieve is \emph{$k$-reachable} for $J,J'$ if $(jj')^k$ sends it
to the maximal sieve, and let $U_k$ be the class of maps all of whose image
and equaliser sieves are $k$-reachable.

\begin{thm}\label{thm:A}
For any topologies $J,J'$ on a small category $S$,
\[
W_J\vee W_{J'}=\bigcup_kU_k,
\]
$U_0$ is the class of isomorphisms, and for $k\ge1$ every map in $U_k$ is a
composite of $2k$ surjections followed by $2k$ injections lying alternately
in $W_{J'}$ and $W_J$. In particular the join is the closure of
$W_J\cup W_{J'}$ under composition.
\end{thm}

\begin{thm}\label{thm:B}
$W_J\vee W_{J'}=W_{J\vee J'}$ if and only if $(jj')^k=j_{J\vee J'}$ on
sieves for some $k$, i.e.\ $(jj')^k=j\vee j'$ as Lawvere--Tierney topologies on
$\Presh(S)$; equivalently, every $(J\vee J')$-covering sieve is $k$-reachable
for one $k$.
\end{thm}

The closure operator of $J\vee J'$ is the transfinite iterate of $jj'$
(Lemma~\ref{lem:joinK}), so Theorem~\ref{thm:B} says that the join of the two
classes is the class of the join topology exactly when the transfinite
iteration stops at a finite stage uniformly. This is the condition
$(j_2j_1)^N=j_1\vee j_2$ that governs joins of nuclei on a Heyting algebra and
their joins in the sameness lattice of fuzzy sets \cite{Laxjoin,Fuzzysame}, now
for the internal frame of a presheaf topos.

For topologies of \emph{ideal type} the criterion becomes a chain
condition. A two-sided ideal $D$ of $S$ is a set of morphisms closed under
composition with arbitrary morphisms on either side; it is \emph{idempotent}
if every member is a composite of two members. Every idempotent ideal $D$
defines a topology $J_D$ whose covering sieves on $c$ are the sieves
containing $D(-,c)=\{a\colon x\to c\text{ in }D\}$ \cite[Thm.~5.1]{Postnikov};
on a finite category every topology is of this form (Kelly--Lawvere
\cite{KL89}, \cite{Marques25}), on an infinite one not necessarily. For
ideals $D,E$ write $DE=\{d\circ e:d\in D,e\in E\}$, again an ideal, and
$(DE)^k=DE\cdots DE$ ($k$ factors of each).

\begin{cor}\label{cor:ideal}
Let $D,E$ be idempotent ideals of $S$. Then $W_{J_D\vee J_E}=W_{J_D}\vee
W_{J_E}$ if and only if the chain $DE\supseteq(DE)^2\supseteq\cdots$
stabilises, i.e.\ $(DE)^m=(DE)^{m+1}$ for some $m$. In that case
$I':=(DE)^m$ is an idempotent ideal with $(ED)^{m+1}=I'$,
$J_D\vee J_E=J_{I'}$, $W_{J_D}\vee W_{J_E}=W_{J_{I'}}$, and every
$J_{I'}$-local isomorphism is a composite of $4m$ maps of
$W_{J_D}\cup W_{J_E}$. In every case $W_{J_D}\vee W_{J_E}$ is the class of
maps locally bijective along some power $(DE)^k$.
\end{cor}

On a finite category the chain stabilises for every pair, which recovers
\cite{Finitejoin} with a proof of a few lines in place of the separated
quotients and plus towers used there; on $\mathbb N$ with the ideals of the
even and of the odd points it is strictly decreasing, which recovers
\cite[Thm.~7.1]{Topjoin}. The criterion becomes concrete for ideals generated
by objects. For a set $A$ of objects, the ideal $D_A$ of morphisms factoring
through $A$ is idempotent, and $W_A=W_{J_{D_A}}$ is the class of maps
bijective at every object of $A$, the class inverted by restriction to $A$,
the sheafification for the essential subtopos
$\Presh(A)\hookrightarrow\Presh(S)$. Proposition~\ref{prop:objects} decides
when $W_A\vee W_B=W_{A\cap B}$: iff every morphism factoring through a long
enough alternating zigzag of objects of $B$ and $A$ factors through an object
of $A\cap B$, on a locally finite poset iff the strictly alternating chains
avoiding $A\cap B$ have bounded size. On $\mathbb N$ with the even and the odd
points the join is the class of maps of towers that are bijective up to a
uniform shift, a proper subclass of the isomorphisms of pro-sets
(Example~\ref{ex:pro}).

The proofs are elementary. The inclusion $U_k\subseteq W_J\vee W_{J'}$ is a
peeling argument: a surjection whose equaliser sieves are $k$-reachable is
divided by the equivalence relation ``equal under $f$ and $J'$-locally
equal'', which is a $J'$-local isomorphism, and the induced map has equaliser
sieves reachable in one step fewer; dually an injection is enlarged to the
$J'$-closure of its image. The failure direction uses, for a covering sieve
$S$ that is not $k$-reachable, the fold map of two copies of a representable
glued along $S$; the coproduct of these maps over all $k$ is inverted by the
join sheafification but lies in no $U_k$.

\subsection*{What is known.}
Samenesses and $W_J$ are from \cite{Same}; that meets of topologies go to
meets of samenesses, the reduction lemma and the $\mathbb N$ counterexample are
from \cite{Topjoin}; the finite case, with the alternating separated
quotients, the plus tower and the retraction of representing bimodules, is
\cite{Finitejoin}. Ideal topologies, the correspondence with idempotent
ideals on finite categories and the local operators $R\mapsto\{a:a\cdot D
\subseteq R\}$ are in \cite{KL89,Marques25,Postnikov}. Local isomorphisms as
locally injective and locally surjective maps are standard
\cite[III.7]{MLM}; joins of nuclei as transfinite iterates are in
\cite{Escardo}. The first version of this paper (archived at DOI
10.5281/zenodo.23467281) proved Corollary~\ref{cor:ideal} for ideal
topologies only, with the stabilising direction by the route of
\cite{Finitejoin}.

\subsection*{What is new.}
Theorem~\ref{thm:A}, the description of the join for arbitrary topologies
and the fact that no cancellations are needed; Theorem~\ref{thm:B}, the
finite-iteration criterion, which decides the question for every pair of
topologies on every small category, whereas \cite{Finitejoin} proves the
positive answer under a finiteness hypothesis; Corollary~\ref{cor:ideal} with
its short proof; Proposition~\ref{prop:objects} and Example~\ref{ex:pro}.

\section{Preliminaries}\label{sec:prelim}

\subsection*{Sieves and closure operators.}
A sieve on $c$ is a set $S$ of morphisms with codomain $c$ closed under
precomposition; for $b\colon x\to c$, $b^*S=\{e:be\in S\}$ is a sieve on $x$.
A Grothendieck topology $J$ assigns to each $c$ a set $J(c)$ of covering
sieves: the maximal sieve is covering, $b^*S$ is covering when $S$ is, and
$S$ is covering when some covering $R$ has $b^*S$ covering for all $b\in R$
\cite[III.2]{MLM}. Covering sieves are closed under finite intersections. The
\emph{closure operator} of $J$ is
\[
j(S)=\{b\colon x\to c:\ b^*S\in J(x)\},
\]
a sieve on $c$ containing $S$; $j$ is monotone, idempotent, preserves finite
intersections, and commutes with pullback: $b^*(jS)=j(b^*S)$. A sieve $S$ on
$c$ is covering iff $\mathrm{id}_c\in j(S)$ iff $j(S)$ is maximal. These
operators are the Lawvere--Tierney topologies on $\Presh(S)$ \cite[V.1]{MLM},
$\Omega(c)$ being the set of sieves on $c$; topologies are ordered by
$J\le J'$ iff $J(c)\subseteq J'(c)$ for all $c$, iff $j\le j'$ pointwise.

\subsection*{Local isomorphisms.}
For a map $f\colon X\to Y$ of presheaves, $y\in Y(c)$ and $x,x'\in X(c)$ put
\[
\Ims_f(y)=\{b\colon x_0\to c:\ y\cdot b\in\im f_{x_0}\},\quad
\Eq(x,x')=\{b\colon x_0\to c:\ x\cdot b=x'\cdot b\}.
\]
Both are sieves, $b^*\Ims_f(y)=\Ims_f(y\cdot b)$ and $b^*\Eq(x,x')=
\Eq(x\cdot b,x'\cdot b)$. By \cite[III.7, Cor.~6]{MLM}, $f\in W_J$ iff
$\Ims_f(y)\in J(c)$ for all $y$ and $\Eq(x,x')\in J(c)$ whenever $f(x)=f(x')$
($f$ is \emph{$J$-locally surjective} and \emph{$J$-locally injective}). In
particular an injection is in $W_J$ iff it is $J$-locally surjective
($J$-dense), and a surjection is in $W_J$ iff it is $J$-locally injective.

\subsection*{Words and reachability.}
Fix topologies $J,J'$ with closure operators $j,j'$ and put $c:=jj'$, a
monotone inflationary operator on sieves commuting with pullback. For an
operator $w$ composed of $j$'s and $j'$'s (a \emph{word}; $w=jj'$ means
$j'$ first), call a sieve $S$ on $c$ \emph{$w$-reachable} if
$\mathrm{id}_c\in w(S)$, equivalently $w(S)$ is maximal; \emph{$k$-reachable}
means $c^k$-reachable.

\begin{lem}\label{lem:reach}
Let $w,w'$ be words and $S,T$ sieves on $c$.
\begin{enumerate}
\item[(i)] $b\in w(S)$ iff $b^*S$ is $w$-reachable. Hence $S$ is
$w$-reachable iff $b^*S$ is $w$-reachable for every $b$ into $c$.
\item[(ii)] If $S$ is $w$-reachable and $b^*T$ is $w'$-reachable for every
$b\in S$, then $T$ is $ww'$-reachable.
\item[(iii)] If $S$ and $T$ are $w$-reachable, so is $S\cap T$; if $S\subseteq
T$ and $S$ is $w$-reachable, so is $T$; a $k$-reachable sieve is
$(k+1)$-reachable; a $J$-covering or $J'$-covering sieve is $1$-reachable.
\end{enumerate}
\end{lem}
\begin{proof}
(i) $b^*(wS)=w(b^*S)$, and $b\in wS$ iff $\mathrm{id}\in b^*(wS)$. (ii) By
(i), $S\subseteq w'(T)$, so $w(S)\subseteq w(w'(T))$ contains the identity.
(iii) Each of $j,j'$ preserves finite intersections, is monotone and is
inflationary; $j(S)\subseteq jj'(S)$ and $j'(S)\subseteq jj'(S)$.
\end{proof}

\subsection*{The join topology.}
Let $K=J\vee J'$ be the least topology above $J$ and $J'$ and $j_K$ its
closure operator. Define the transfinite iterates $c^0(S)=S$,
$c^{\alpha+1}(S)=c(c^\alpha(S))$ and $c^\lambda(S)=\bigcup_{\alpha<\lambda}
c^\alpha(S)$ for limit $\lambda$; they increase with $\alpha$, commute with
pullback, and $c^\alpha(c^\beta(S))=c^{\beta+\alpha}(S)$.

\begin{lem}\label{lem:joinK}
$j_K(S)=\bigcup_\alpha c^\alpha(S)$; a sieve is $K$-covering iff it is
$c^\alpha$-reachable for some ordinal $\alpha$. (The sieves on $c$ form a
set, so the increasing iterates $c^\alpha(S)$ are constant from some ordinal
on, and the union is attained.)
\end{lem}
\begin{proof}
Since $j,j'\le j_K$ and $j_K$ is monotone and idempotent, $c^\alpha(S)\subseteq
j_K(S)$ for all $\alpha$ by induction (unions at limits). Let $T(c)$ be the set
of sieves $S$ on $c$ with $\mathrm{id}_c\in c^\alpha(S)$ for some $\alpha$. It
contains $J(c)$ and $J'(c)$ by Lemma~\ref{lem:reach}(iii), the maximal sieve,
and is stable under pullback since $b^*(c^\alpha S)=c^\alpha(b^*S)$. For
transitivity let $S\in T(c)$ and $b^*S'\in T$ for all $b\in S$, say
$\mathrm{id}\in c^{\alpha_b}(b^*S')$; with $\beta=\sup_b\alpha_b$ we get
$b\in c^\beta(S')$ for all $b\in S$, so $S\subseteq c^\beta(S')$ and
$\mathrm{id}\in c^\alpha(S)\subseteq c^{\beta+\alpha}(S')$. Thus $T$ is a
topology above $J,J'$, hence $T\supseteq K$, and $T\subseteq K$ by the first
sentence. Finally $b\in j_K(S)$ iff $b^*S\in K$ iff $\mathrm{id}\in
c^\alpha(b^*S)=b^*(c^\alpha S)$ for some $\alpha$ iff $b\in\bigcup_\alpha
c^\alpha(S)$.
\end{proof}

\section{Uniform locality and the join}\label{sec:uniform}

Fix topologies $J,J'$, $K=J\vee J'$, $c=jj'$, and write
$V:=W_J\vee W_{J'}$.

\begin{defi}
For $k\ge0$ let $U_k$ be the class of maps $f\colon X\to Y$ such that every
image sieve $\Ims_f(y)$ and every equaliser sieve $\Eq(x,x')$ with $f(x)=f(x')$
is $k$-reachable. Put $U=\bigcup_kU_k$.
\end{defi}

$U_0$ is the class of isomorphisms (maximal image sieves mean surjectivity,
maximal equaliser sieves injectivity). A map in $U_k$ is a $K$-local
isomorphism \emph{uniformly}: the ordinal of Lemma~\ref{lem:joinK} needed for
its image and equaliser sieves is the finite number $k$, independent of the
element.

\begin{lem}\label{lem:U}
$U_k\subseteq U_{k+1}$, $W_J\cup W_{J'}\subseteq U_1$, $U\subseteq W_K$, and
$U$ is a sameness. Hence $V\subseteq U\subseteq W_K$.
\end{lem}
\begin{proof}
The first three assertions are Lemma~\ref{lem:reach}(iii) and
Lemma~\ref{lem:joinK}. Identities lie in $U_0$. Let $h\colon X\to Y$ and
$g\colon Y\to Z$.

\emph{Composition.} Let $h\in U_k$, $g\in U_l$. For $z\in Z(c)$ the sieve
$S=\Ims_g(z)$ is $l$-reachable; for $b\in S$ write $z\cdot b=g(y_b)$; then
$b^*\Ims_{gh}(z)=\Ims_{gh}(z\cdot b)\supseteq\Ims_h(y_b)$ is $k$-reachable, so
$\Ims_{gh}(z)$ is $(l+k)$-reachable by Lemma~\ref{lem:reach}(ii). If
$gh(x)=gh(x')$, the sieve $S=\Eq(h(x),h(x'))$ is $l$-reachable; for $b\in S$,
$h(x\cdot b)=h(x'\cdot b)$, so $b^*\Eq(x,x')=\Eq(x\cdot b,x'\cdot b)$ is
$k$-reachable; hence $\Eq(x,x')$ is $(l+k)$-reachable.

\emph{Cancellation of $g$.} Let $gh\in U_k$, $g\in U_k$. If $h(x)=h(x')$
then $gh(x)=gh(x')$ and $\Eq(x,x')$ is $k$-reachable. For $y\in Y(c)$ the
sieve $S=\Ims_{gh}(g(y))$ is $k$-reachable; for $b\in S$ write
$g(y)\cdot b=gh(x_b)$, so $g(y\cdot b)=g(h(x_b))$ and $T_b=\Eq(y\cdot b,h(x_b))$
is $k$-reachable; for $e\in T_b$, $y\cdot be=h(x_b\cdot e)\in\im h$, so
$b^*\Ims_h(y)\supseteq T_b$. By Lemma~\ref{lem:reach}(ii), $\Ims_h(y)$ is
$2k$-reachable.

\emph{Cancellation of $h$.} Let $gh\in U_k$, $h\in U_k$. Then
$\Ims_g(z)\supseteq\Ims_{gh}(z)$ is $k$-reachable. If $g(y)=g(y')$, the sieve
$S=\Ims_h(y)\cap\Ims_h(y')$ is $k$-reachable; for $b\in S$ write $y\cdot b=
h(x_b)$, $y'\cdot b=h(x'_b)$; then $gh(x_b)=g(y)\cdot b=g(y')\cdot b=gh(x'_b)$,
so $\Eq(x_b,x'_b)$ is $k$-reachable, and for $e\in\Eq(x_b,x'_b)$,
$y\cdot be=h(x_b\cdot e)=h(x'_b\cdot e)=y'\cdot be$; so $b^*\Eq(y,y')\supseteq
\Eq(x_b,x'_b)$ and $\Eq(y,y')$ is $2k$-reachable.
\end{proof}

\begin{lem}[peeling]\label{lem:peel}
Let $w=j_1j_2\cdots j_n$ be a word, each $j_i\in\{j,j'\}$ with topology
$J_i$.
\begin{enumerate}
\item[(a)] If $f\colon X\to Y$ is surjective and $\Eq(x,x')$ is $w$-reachable
whenever $f(x)=f(x')$, then $f=f_1\circ\cdots\circ f_n$ with $f_i$ a
surjection in $W_{J_i}$.
\item[(b)] If $f\colon X\to Y$ is injective and every $\Ims_f(y)$ is
$w$-reachable, then $f=f_1\circ\cdots\circ f_n$ with $f_i$ an injection in
$W_{J_i}$.
\end{enumerate}
\end{lem}
\begin{proof}
Induction on $n$, proving that $f$ is a composite of $n$ such maps followed
by an isomorphism, which for $n\ge1$ can be absorbed into the last factor;
for $n=0$ the hypothesis says that $f$ is bijective.
(a) Let $R$ be the relation on $X(c)$: $x\,R\,x'$ iff $f(x)=f(x')$ and
$\Eq(x,x')\in J_n(c)$. It is an equivalence relation, since covering sieves
contain the maximal sieve and are closed under intersection, and a
congruence, since $b^*\Eq(x,x')=\Eq(x\cdot b,x'\cdot b)$ and covering sieves
pull back. The quotient $q\colon X\twoheadrightarrow X/R$ is $J_n$-locally
injective by construction, hence in $W_{J_n}$, and $f$ factors as
$\bar f\circ q$ with $\bar f$ surjective. If $\bar f[x]=\bar f[x']$ then
\[
\Eq([x],[x'])=\{b:\ x\cdot b\,R\,x'\cdot b\}=\{b:\ b^*\Eq(x,x')\in J_n\}
=j_n(\Eq(x,x')),
\]
and since $\mathrm{id}\in j_1\cdots j_{n-1}(j_n\Eq(x,x'))$ this sieve is
$j_1\cdots j_{n-1}$-reachable. Induction applied to $\bar f$ gives the
factorisation. (b) Regard $X$ as a subpresheaf of $Y$ and let
$Z(c)=\{y\in Y(c):\Ims_f(y)\in J_n(c)\}$, the $J_n$-closure of $X$ in $Y$; it is
a subpresheaf containing $X$, the inclusion $X\hookrightarrow Z$ is
$J_n$-dense, hence in $W_{J_n}$, and for $y\in Y(c)$,
\[
\Ims_Z(y)=\{b:\ y\cdot b\in Z\}=\{b:\ b^*\Ims_f(y)\in J_n\}=j_n(\Ims_f(y)),
\]
which is $j_1\cdots j_{n-1}$-reachable. Induction applied to
$Z\hookrightarrow Y$ gives the factorisation.
\end{proof}

\begin{proof}[Proof of Theorem~\ref{thm:A}]
Let $C$ be the closure of $W_J\cup W_{J'}$ under composition. Then
$C\subseteq V\subseteq U$ by Lemma~\ref{lem:U}. Let $f\in U_k$ and factor
$f=m\circ e$ with $e\colon X\twoheadrightarrow\im f$ and $m\colon\im f
\hookrightarrow Y$. The equaliser sieves of $e$ are those of $f$, and the
image sieves of $m$ are those of $f$; both are $c^k$-reachable, and $c^k=
(jj')^k$ is a word of length $2k$ ending with $j'$. By Lemma~\ref{lem:peel},
$e$ is a composite of $2k$ surjections and $m$ a composite of $2k$
injections, alternately in $W_{J'}$ and $W_J$, beginning with $W_{J'}$. So
$U\subseteq C$, and $C=V=U$.
\end{proof}

\begin{figure}[ht]
\[
\begin{tikzcd}[column sep=normal,row sep=large]
X\arrow[r,two heads,"W_{J'}"]\arrow[rrr,bend left=25,"e"] & X/R_1\arrow[r,two heads,"W_J"] & \cdots\arrow[r,two heads,"W_J"] & \im f\arrow[d,hook,"W_{J'}"]\arrow[ddlll,bend left=35,"m"{near start}]\\
& & & Z_1\arrow[d,hook,"W_J"]\\
Y & \cdots\arrow[l,hook',"W_J"'] & & Z_2\arrow[ll,hook',"W_{J'}"']
\end{tikzcd}
\]
\caption{Theorem~\ref{thm:A}: a map in $U_k$ is $2k$ quotients by
``equal under $f$ and locally equal'' followed by $2k$ alternating closures
of its image, each step a local isomorphism for one of the two topologies.}
\end{figure}

\begin{rem}
In general the join of two samenesses requires the two-out-of-three
saturation of the compositional closure \cite{Same}; here no cancellation is
needed. The number of steps is controlled in both directions: a composite of
$n$ maps of $W_J\cup W_{J'}$ lies in $U_n$ by Lemma~\ref{lem:U}, and a map of
$U_k$ is a composite of $4k$ such maps.
\end{rem}

\section{The criterion}\label{sec:criterion}

\begin{prop}\label{prop:witness}
Let $S$ be a sieve on $c$, let $X_S$ be the presheaf obtained from two copies
$y(c)^{(1)},y(c)^{(2)}$ of the representable $y(c)=S(-,c)$ by identifying
$b^{(1)}$ with $b^{(2)}$ for every $b\in S$ (the pushout of $y(c)\leftarrow S
\rightarrow y(c)$), and let $f_S\colon X_S\to y(c)$ be the fold map. Then
$f_S\in W_K$ iff $S$ is $K$-covering, and $f_S\in U_k$ iff $S$ is
$k$-reachable.
\end{prop}
\begin{proof}
The identification is compatible with restriction because $S$ is a sieve,
so $X_S$ is a presheaf and $f_S$ a surjective map of presheaves; its image
sieves are maximal. Two distinct elements with the same image are $b^{(1)}$,
$b^{(2)}$ for some $b\notin S$, and $\Eq(b^{(1)},b^{(2)})=\{e:be\in S\}=b^*S$.
So $f_S\in W_K$ iff $b^*S\in K$ for all $b\notin S$; if $S$ is maximal $f_S$
is an isomorphism, otherwise $\mathrm{id}_c\notin S$ and the condition
implies $S\in K$, while $S\in K$ implies $b^*S\in K$ for all $b$. Likewise
$f_S\in U_k$ iff $b^*S$ is $k$-reachable for all $b\notin S$ iff $b\in c^k(S)$
for all $b\notin S$ (Lemma~\ref{lem:reach}(i)) iff $c^k(S)$ is maximal, as
$c^k(S)\supseteq S$.
\end{proof}

\begin{figure}[ht]
\[
\begin{tikzcd}[column sep=large]
S\arrow[r,hook]\arrow[d,hook] & y(c)^{(2)}\arrow[d]\arrow[ddr,bend left=15,"\mathrm{id}"]\\
y(c)^{(1)}\arrow[r]\arrow[drr,bend right=15,"\mathrm{id}"'] & X_S\arrow[dr,"f_S"]\\
& & y(c)
\end{tikzcd}
\]
\caption{The witness of Proposition~\ref{prop:witness}: two copies of the
representable glued along the sieve $S$; the fold map identifies the two
identities, whose equaliser sieve is $S$.}
\end{figure}

\begin{proof}[Proof of Theorem~\ref{thm:B}]
The two formulations of the condition agree: $(jj')^k\le j_K$ always, and if
every $K$-covering sieve is $k$-reachable then $b\in j_K(S)$ gives $b^*S\in K$,
so $\mathrm{id}\in c^k(b^*S)=b^*c^k(S)$ and $b\in c^k(S)$; conversely
$c^k=j_K$ makes every covering sieve $k$-reachable. If the condition holds,
a map in $W_K$ has $K$-covering, hence $k$-reachable, image and equaliser
sieves, so $W_K\subseteq U_k\subseteq V$, and $V\subseteq W_K$ always. If it
fails, choose for each $k$ an object $c_k$ and a $K$-covering sieve $S_k$ on
$c_k$ that is not $k$-reachable, and let $F=\coprod_kf_{S_k}$. Then $F\in W_K$,
since image and equaliser sieves of a coproduct of maps are those of the
components (elements of a coproduct of presheaves restrict within their
summand), and $F\notin U_k$ for every $k$ because its component $f_{S_k}$ is
not in $U_k$ (Proposition~\ref{prop:witness}). So $F\in W_K\setminus U=
W_K\setminus V$ by Theorem~\ref{thm:A}.
\end{proof}

\begin{rem}
Lemma~\ref{lem:joinK} and Theorem~\ref{thm:B} place the question in the
setting of \cite{Laxjoin,Fuzzysame}: there the join $j_1\vee j_2$ of two
nuclei on a Heyting algebra or frame $L$ is the transfinite iterate of
$j_2j_1$ \cite{Escardo}, and the identity $(j_2j_1)^N=j_1\vee j_2$ for some
$N$, the finite-iteration axiom of the lax logic of the join in
\cite{Laxjoin}, is exactly the condition under which the join of the two
reflections of $L$-fuzzy sets is generated at the level of inverted maps
\cite[Cor.~4.3]{Fuzzysame}. Theorem~\ref{thm:B} is the same statement for
the internal frame $\Omega$ of $\Presh(S)$ with the Lawvere--Tierney
topologies $j,j'$: the join of two subtopos classes is the class of the meet
of the subtoposes iff the alternation of the two local operators reaches the
join uniformly in finitely many steps. The condition is uniform in the
sieve; it should be distinguished from the pointwise condition that the
iteration terminates at each element, which suffices in $\Same(L)$
\cite{Fuzzysame} and which, for the relativized necessities
$\Box_C,\Box_{C'}$ on the frame of up-sets of a poset, is the absence of
infinite alternating chains characterising $\Box_C\vee\Box_{C'}=
\Box_{C\cap C'}$ \cite{Relnec}; a disjoint union of finite alternating chains
of unbounded length satisfies the pointwise condition but not the uniform
one.
\end{rem}

\section{Topologies of ideal type}\label{sec:ideal}

For an ideal $A$ the sets $A(-,c)$ are sieves and $A(-,c)\cdot b\subseteq
A(-,x)$ for $b\colon x\to c$.

\begin{lem}\label{lem:ideal}
Let $D,E$ be idempotent ideals.
\begin{enumerate}
\item[(i)] $J_D$ is a topology, $D(-,c)$ is its least covering sieve on $c$,
and $J_D\le J_E$ iff $D\supseteq E$.
\item[(ii)] The closure operator of $J_D$ is $j_D(S)=\{b:\ b\cdot D\subseteq
S\}$, where $b\cdot D=\{b\circ d:d\in D(-,\operatorname{dom}b)\}$; hence
$(j_Dj_E)^k(S)=\{b:\ b\cdot(DE)^k\subseteq S\}$, and $S$ is $k$-reachable for
$J_D,J_E$ iff $S\supseteq(DE)^k(-,c)$. A map is in $U_k$ iff it is
\emph{locally bijective along $(DE)^k$}: $f(x)=f(x')$ implies $x\cdot a=
x'\cdot a$ for all $a\in(DE)^k(-,c)$, and $y\cdot a\in\im f$ for all $y\in Y(c)$
and $a\in(DE)^k(-,c)$.
\item[(iii)] For every $k$ and $c$ the sieve $(DE)^k(-,c)$ is
$(J_D\vee J_E)$-covering.
\item[(iv)] If $(DE)^m=(DE)^{m+1}=:I'$ then $I'$ is idempotent,
$I'\subseteq D\cap E$, $DI'=I'D=EI'=I'E=I'$, $(ED)^{m+1}=I'$, every
alternating product of at least $2(m+1)$ factors $D,E$, in either order,
equals $I'$, and $J_D\vee J_E=J_{I'}$. (The chain $(ED)^k$ need not have
stabilised at index $m$: Example~\ref{ex:ED}.)
\end{enumerate}
\end{lem}
\begin{proof}
(i) is \cite[Thm.~5.1]{Postnikov}: the maximal sieve contains $D(-,c)$;
pullback of a sieve containing $D(-,c)$ along $b\colon x\to c$ contains
$D(-,x)$ because $D$ is a left ideal; and if $R\supseteq D(-,c)$ and each
$a^*R'$ with $a\in R$ contains $D(-,\operatorname{dom}a)$, then for
$d\in D(-,c)$, $d=d_1d_2$ with $d_i\in D$ by idempotency, $d_1\in R$ and
$d_2\in d_1^*R'$, so $d\in R'$. (ii) $b\in j_D(S)$ iff $b^*S\supseteq
D(-,\operatorname{dom}b)$ iff $b\cdot D\subseteq S$; then $b\in j_Dj_E(S)$ iff
$(b\circ d)\cdot E\subseteq S$ for all $d$ iff $b\cdot DE\subseteq S$, and
inductively for $(DE)^k$; $\mathrm{id}_c\cdot(DE)^k=(DE)^k(-,c)$. For a map
$f$, $\Ims_f(y)\supseteq(DE)^k(-,c)$ says $y\cdot a\in\im f$ for $a\in(DE)^k
(-,c)$, and $\Eq(x,x')\supseteq(DE)^k(-,c)$ says $x\cdot a=x'\cdot a$. (iii)
is Lemma~\ref{lem:joinK}: $(DE)^k(-,c)$ is $k$-reachable. (iv) $I'I'=(DE)^{2m}
=I'$; $I'\subseteq DE\subseteq D\cap E$; $DI'=DD\,E(DE)^{m-1}=(DE)^m=I'$ as
$DD=D$, and similarly $I'E=I'$; $I'D\supseteq I'I'=I'$ as $I'\subseteq D$, and
$I'D\subseteq I'$ as $I'$ is a right ideal, so $I'D=I'$, and likewise $EI'=I'$;
then $(ED)^{m+1}=E(DE)^mD=EI'D=I'$. An alternating product of at least
$2(m+1)$ factors contains $(DE)^{m+1}$ or $(ED)^{m+1}$ as a block of
consecutive factors; the block equals $I'$, and multiplying $I'$ by $D$ or $E$
on either side returns $I'$. For the join, $J_{I'}\ge J_D,J_E$ since
$I'\subseteq D,E$, and any topology above $J_D$ and $J_E$ has
$I'(-,c)=(DE)^m(-,c)$ covering by (iii), hence is above $J_{I'}$.
\end{proof}

\begin{proof}[Proof of Corollary~\ref{cor:ideal}]
By Lemma~\ref{lem:ideal}(ii) and Theorem~\ref{thm:A}, $W_{J_D}\vee W_{J_E}$
is the class of maps locally bijective along some $(DE)^k$. If
$(DE)^m=(DE)^{m+1}=I'$, then $K=J_{I'}$ by Lemma~\ref{lem:ideal}(iv), a
$K$-covering sieve contains $I'(-,c)=(DE)^m(-,c)$ and is therefore
$m$-reachable; Theorem~\ref{thm:B} gives $V=W_K=W_{J_{I'}}$, and
Theorem~\ref{thm:A} gives the factorisation into $4m$ maps, since
$W_{J_{I'}}=U_m$. If the chain does not stabilise, then $(DE)^k\ne(DE)^{k+1}$
for every $k$ (once two consecutive terms agree, all later ones do), so there
are $c_k$ with $(DE)^k(-,c_k)\not\subseteq(DE)^{k+1}(-,c_k)$; the sieve
$S_k=(DE)^{k+1}(-,c_k)$ is $K$-covering by Lemma~\ref{lem:ideal}(iii) and not
$k$-reachable by Lemma~\ref{lem:ideal}(ii), so Theorem~\ref{thm:B} gives
$V\ne W_K$. The witnesses of Proposition~\ref{prop:witness} are here two
copies of $y(c_k)$ glued along $(DE)^{k+1}(-,c_k)$, whose fold map is locally
bijective along $(DE)^{k+1}$ but identifies the two identities, which
restrict along a morphism of $(DE)^k\setminus(DE)^{k+1}$ into $c_k$ to
distinct elements.
\end{proof}

\begin{rem}
The first version of this paper proved the stabilising direction by the
route of \cite{Finitejoin}: $2r$ alternating separated quotients, $r=m+1$,
followed by $4r$ alternating plus constructions, with a retraction of the
representing bimodule $I'\otimes I'$ onto the alternating tensor power.
Theorem~\ref{thm:A} replaces this by the single observation
$W_{J_{I'}}=U_m\subseteq V$ and improves the bound to $4m$ maps. The plus
tower is nevertheless of independent interest, as it need not stabilise even
when the chain does \cite[\S7]{Finitejoin}.
\end{rem}

\begin{ex}\label{ex:ED}
The two chains $(DE)^k$ and $(ED)^k$ need not stabilise at the same index.
Let $S$ be the category with two objects $0,1$ and one non-identity arrow
$a\colon0\to1$, and let $D=\{\mathrm{id}_0,a\}$, $E=\{\mathrm{id}_1,a\}$ be the
ideals of arrows factoring through $0$, respectively through $1$; both are
idempotent. Then $DE=\emptyset$ (no arrow of $E$ has codomain $0$), while
$ED=\{a\}$. So $(DE)^1=(DE)^2=\emptyset=I'$, $m=1$, but $(ED)^1=\{a\}\ne I'$,
and only $(ED)^2=\emptyset=I'$. The join $J_D\vee J_E=J_\emptyset$ is the
degenerate topology, as the empty sieve on $0$ is $J_E$-covering and the
sieve $\{a\}$ on $1$ is $J_D$-covering.
\end{ex}

\begin{ex}\label{ex:N}
On a finite category every chain of ideals stabilises, so
Corollary~\ref{cor:ideal} contains \cite[Thm.~1.1]{Finitejoin}. On
$S=(\mathbb N,\le)$ with $D$ and $E$ the ideals of arrows $n\to m$ passing
through an even, respectively an odd, point, $(DE)^k$ consists of the arrows
$n\to m$ admitting $k$ alternating even and odd points between $n$ and $m$,
a strictly decreasing chain with empty intersection; the join topology is
degenerate and $W_K$ is everything, while $V$ is not \cite[Thm.~7.1]{Topjoin}.
Example~\ref{ex:pro} describes $V$.
\end{ex}

\begin{ex}
For a monoid $M$, viewed as a category with one object, ideals are the
two-sided ideals of $M$ and the criterion is the stabilisation of the
alternating products $DE\supseteq DEDE\supseteq\cdots$ in the quantale of
ideals of $M$. Every finite monoid satisfies it, with $m$ at most the number
of elements. For infinite monoids the chain can decrease forever, and then
the join of the two classes of local isomorphisms of $M$-sets is the class
of maps locally bijective along some $(DE)^k$, strictly smaller than the
class for the join topology.
\end{ex}

\section{Ideals generated by objects}\label{sec:objects}

For a set $A$ of objects of $S$ let $D_A$ be the set of morphisms factoring
through an object of $A$. It is an ideal, idempotent since $f=g\circ h$ with
$h\colon x\to a$, $g\colon a\to y$ has $g,h\in D_A$; on a finite category
every idempotent ideal is of this form after Cauchy completion
\cite[\S2]{Finitejoin}, but not in general: on $(\mathbb Q,\le)$ the ideal of
strict inequalities is idempotent and contains no identity. Write
$W_A:=W_{J_{D_A}}$.

\begin{prop}\label{prop:objects}
Let $A,B$ be sets of objects of $S$.
\begin{enumerate}
\item[(a)] $W_A$ is the class of maps $f$ with $f_a$ bijective for every
$a\in A$; $(D_AD_B)^k$ is the set of morphisms factoring through a chain
$b_1\to a_1\to b_2\to\cdots\to b_k\to a_k$ with $b_i\in B$, $a_i\in A$, and
contains $D_{A\cap B}$. The join $W_A\vee W_B$ is the class of maps $f$ such
that, for some $k$, $f$ is bijective along every such chain: if $f(x)=f(x')$
then $x\cdot g=x'\cdot g$, and every $y\cdot g$ lies in the image, for every
morphism $g$ factoring through a chain of length $k$.
\item[(b)] $W_A\vee W_B=W_{A\cap B}$ if and only if for some $N$ every
morphism factoring through such a chain of length $N$ factors through an
object of $A\cap B$, i.e.\ $(D_AD_B)^N\subseteq D_{A\cap B}$.
\item[(c)] If $S=P$ is a poset whose intervals are finite, the condition in
(b) holds if and only if the strictly alternating chains
$b_1<a_1<\cdots<b_k<a_k$ ($b_i\in B$, $a_i\in A$) contained in intervals
disjoint from $A\cap B$ have a bounded number of elements; if they do not,
the chain $(D_AD_B)^k$ does not stabilise and $W_A\vee W_B\subsetneq
W_{J_{D_A}\vee J_{D_B}}$.
\end{enumerate}
\end{prop}
\begin{proof}
(a) If $f_a$ is bijective for $a\in A$, then for $x,x'\in X(c)$ with
$f(x)=f(x')$ and $g=g_1g_2\in D_A(-,c)$ with $g_1\colon a\to c$, $a\in A$, we
get $f_a(x\cdot g_1)=f_a(x'\cdot g_1)$, so $x\cdot g_1=x'\cdot g_1$ and
$x\cdot g=x'\cdot g$; and for $y\in Y(c)$, $y\cdot g_1=f_a(z)$ for some $z$,
so $y\cdot g=f(z\cdot g_2)$. Thus $f$ is locally bijective along $D_A$, and
$f\in W_A$ by Lemma~\ref{lem:ideal}(ii). Conversely $\mathrm{id}_a\in
D_A(-,a)$ for $a\in A$, so local bijectivity along $D_A$ gives $f_a$
bijective. A morphism of $D_AD_B$ is $d\circ e$ with $e$ factoring through
some $b\in B$ and $d$ through some $a\in A$, i.e.\ a morphism factoring
through a chain $b\to a$; iterating gives the description of $(D_AD_B)^k$,
and a morphism through $c\in A\cap B$ factors through the chain
$c\to c\to\cdots\to c$ of identities. The description of the join is
Corollary~\ref{cor:ideal}.
(b) If $(D_AD_B)^N\subseteq D_{A\cap B}$ then
$(D_AD_B)^N\subseteq D_{A\cap B}\subseteq(D_AD_B)^{N+1}\subseteq(D_AD_B)^N$,
so the chain stabilises at $I'=D_{A\cap B}$ and Corollary~\ref{cor:ideal}
gives $W_A\vee W_B=W_{J_{I'}}=W_{A\cap B}$. Conversely, $D_{A\cap B}\subseteq
D_A\cap D_B$ gives $J_{D_{A\cap B}}\ge K:=J_{D_A}\vee J_{D_B}$ by
Lemma~\ref{lem:ideal}(i), hence $W_{A\cap B}\supseteq W_K\supseteq
W_A\vee W_B$. If $W_A\vee W_B=W_{A\cap B}$, then $W_A\vee W_B=W_K$, the chain
stabilises at some $I'=(D_AD_B)^m$ by Corollary~\ref{cor:ideal}, and
$W_{J_{I'}}=W_K=W_{J_{D_{A\cap B}}}$. A topology is determined by the class
of maps inverted by its sheafification, since the sheaves are the objects
local with respect to that class (a reflective subcategory is recovered from
the maps its reflector inverts) and subtoposes of $\Presh(S)$ correspond
bijectively to topologies \cite[Chap.~V]{MLM}; so $J_{I'}=J_{D_{A\cap B}}$,
$I'=D_{A\cap B}$ by Lemma~\ref{lem:ideal}(i), and $N=m$ works.
(c) In a poset, $(n\le m)\in(D_AD_B)^k$ iff there are $n\le b_1\le a_1\le
\cdots\le b_k\le a_k\le m$ with $b_i\in B$, $a_i\in A$; if the interval
$[n,m]$ is disjoint from $A\cap B$, an equality $b_i=a_i$ or $a_i=b_{i+1}$
would produce a point of $A\cap B$, so the chain is strict. Hence
$(D_AD_B)^k\setminus D_{A\cap B}$ consists of the pairs $n\le m$ with $[n,m]$
disjoint from $A\cap B$ and containing a strictly alternating chain with
$2k$ elements. If these chains have at most $L$ elements, then
$(D_AD_B)^k=D_{A\cap B}$ for $2k>L$. If their sizes are unbounded,
$(D_AD_B)^k\ne D_{A\cap B}$ for every $k$, while $\bigcap_k(D_AD_B)^k=
D_{A\cap B}$ because a chain with more elements than the finite interval
$[n,m]$ is not strict; a descending chain whose members all differ from its
intersection does not stabilise, and Corollary~\ref{cor:ideal} gives the
strict inclusion.
\end{proof}

\begin{ex}
On $P=(\mathbb N,\le)$ let $A$ be the even numbers and $B$ the odd numbers
together with the multiples of $4$, so $A\cap B=4\mathbb N$. An interval
disjoint from $4\mathbb N$ lies in some $[4j+1,4j+3]$, whose longest strictly
alternating chain is $4j+1<4j+2<4j+3$; so $L=3$, $(D_AD_B)^2=D_{4\mathbb N}$,
and a map of presheaves on $\mathbb N$ lies in the join of ``bijective at the
even points'' and ``bijective at the odd points and at the multiples of $4$''
if and only if it is bijective at the multiples of $4$.
\end{ex}

\begin{ex}\label{ex:pro}
On $P=(\mathbb N,\le)$ let $A$ be the even and $B$ the odd numbers. A presheaf
is a tower $\cdots\to X(2)\to X(1)\to X(0)$ of sets, and
$(D_AD_B)^k(-,c)$ consists of the $n\le c$ with $c-n\ge2k-1$, or $\ge2k$
when $n$ is even. By Proposition~\ref{prop:objects}(a), $W_A\vee W_B$ is the
class of maps of towers $f\colon X\to Y$ that are \emph{bijective up to a
shift}: for some $s$, $f_c(x)=f_c(x')$ implies $x|_{c-s}=x'|_{c-s}$, and every
$y\in Y(c)$ has $y|_{c-s}\in\im f_{c-s}$ (for $c\ge s$; for $c<s$ the
conditions are void). Such a map becomes an isomorphism of pro-objects: it
has the inverse up to shift $g_n\colon Y(n+2s)\to X(n)$, $y\mapsto x|_n$ for
any $x\in X(n+s)$ with $f(x)=y|_{n+s}$, which is well defined and compatible
with restriction. The converse fails: a pro-isomorphism of towers only needs,
for each $n$, \emph{some} later stage $m(n)$ at which the inverse is defined,
and the join requires the uniform bound $m(n)=n+2s$. For instance let
$X(n)=\{0,1\}$ for all $n$, with the restriction $X(n+1)\to X(n)$ constant
with value $0$ when $n$ is a power of $2$ and the identity otherwise, and let
$f\colon X\to1$ be the map to the terminal tower. The section $1\to X$
picking $0$ is a pro-inverse, since for every $n$ the composite restriction
from some later stage to $n$ is constant; but on an interval
$[n,n+s]$ containing no power of $2$ the two elements of $X(n+s)$ are
identified by $f$ and remain distinct at $n$, so $f$ is bijective up to no
uniform shift and $f\notin W_A\vee W_B$ (although $f\in W_K$, as $W_K$ is the
class of all maps). So in $\Same(\Presh(\mathbb N))$ the join of ``bijective
at the even points'' and ``bijective at the odd points'' is a uniform
version of pro-isomorphism, strictly between the two classes and the
pro-isomorphisms. The obstruction $\lim_{\mathbb N^{\mathrm{op}}}$ used in
\cite[Thm.~7.1]{Topjoin} inverts every pro-isomorphism, and more (the
projection $\{m\ge n\}\times\{0,1\}\to\{m\ge n\}$ has empty limits on both
sides and is no pro-isomorphism), so it separates the join from $W_K$
without describing it: $W_A\vee W_B\subsetneq W_{\text{pro-iso}}\subsetneq
W_{\lim}\subseteq W_K$.
\end{ex}

\begin{rem}
For a set $A$ of objects, the restriction functor along the full inclusion
$i\colon A\hookrightarrow S$ preserves limits and has a fully faithful right
adjoint, the right Kan extension $\mathrm{Ran}_i$ (and a left adjoint, the
left Kan extension). Its essential image is a reflective subcategory whose
reflector inverts exactly $W_A$, hence it is the subcategory of
$J_{D_A}$-sheaves: $\Sh(S,J_{D_A})\simeq\Presh(A)$, an essential
subtopos of $\Presh(S)$ in the sense of \cite{KL89}. The meet of the
subtoposes $\Presh(A)$ and $\Presh(B)$ is the subtopos of
$(J_{D_A}\vee J_{D_B})$-sheaves. Proposition~\ref{prop:objects}(b) says that
$W_A\vee W_B=W_{A\cap B}$ iff this meet is equivalent to $\Presh(A\cap B)$
\emph{and} is reached by a bounded alternation; the first condition alone
does not suffice: for the even and the odd numbers on $\mathbb N$ the meet is
equivalent to the degenerate subtopos $\Presh(\emptyset)$, yet the join of
the samenesses is not $W_\emptyset$ but the class of maps bijective up to a
uniform shift.
\end{rem}

\section{Remarks and questions}\label{sec:rem}

\begin{rem}
The proofs of Section~\ref{sec:uniform} use only the closure operators of
the two topologies on subobjects, image factorisations, quotients by
equivalence relations and the description of local isomorphisms by local
injectivity and surjectivity. In an arbitrary topos with Lawvere--Tierney
topologies $j,j'$ one can define $U_k$ by requiring that $Y$ be the
$2k$-fold alternating closure of $\im f$ and that the kernel pair of $f$ be
contained in the $2k$-fold alternating closure of the diagonal of $X$; the
peeling argument then goes through with closures of subobjects, and we
expect Theorems~\ref{thm:A} and~\ref{thm:B} to hold in this generality, with
Lemma~\ref{lem:U} proved in the internal language. We have not carried this
out.
\end{rem}

\begin{que}
Theorem~\ref{thm:B} reduces join preservation to the identity
$(jj')^k=j\vee j'$ on $\Omega$. Which pairs of topologies on a given infinite
category satisfy it? For ideal topologies this is the chain condition of
Corollary~\ref{cor:ideal} and for object-generated ideals the zigzag condition
of Proposition~\ref{prop:objects}; a structural description for topologies
that are not of ideal type, for instance in terms of the two subtoposes,
is open.
\end{que}

\section*{Funding}
No funding was received for this work.

\section*{Declarations}
The author has no competing interests to declare.

\begin{thebibliography}{9}
\bibitem{Escardo} M.~H.~Escard\'o, \emph{Joins in the frame of nuclei}, Appl.
Categ. Structures \textbf{11} (2003), 117--124.
\bibitem{KL89} G.~M.~Kelly, F.~W.~Lawvere, \emph{On the complete lattice of
essential localizations}, Bull. Soc. Math. Belg. S\'er. A \textbf{41} (1989),
289--319.
\bibitem{MLM} S.~Mac Lane, I.~Moerdijk, \emph{Sheaves in Geometry and Logic},
Springer, 1992.
\bibitem{Marques25} J.~Marqu\`es, \emph{A criterion for categories on which
every Grothendieck topology is rigid}, Appl. Categ. Structures \textbf{33}
(2025), art.\ 37.
\bibitem{Same} R.~Shibusawa, \emph{The sameness lattice of a category:
weak-equivalence structures, Galois functoriality, and comparison intervals},
preprint (2026), archived at DOI 10.5281/zenodo.22823772.
\bibitem{Postnikov} R.~Shibusawa, \emph{Postnikov truncations as joins of sheaf
and homotopy equivalences}, preprint (2026), archived at DOI 10.5281/zenodo.22956181.
\bibitem{Topjoin} R.~Shibusawa, \emph{Meets and joins of Grothendieck
topologies in the sameness lattice}, preprint (2026), archived at DOI
10.5281/zenodo.22978446.
\bibitem{Finitejoin} R.~Shibusawa, \emph{Grothendieck topologies on a finite
category form a sublattice of the lattice of two-out-of-three classes},
preprint (2026), archived at DOI 10.5281/zenodo.23204943.
\bibitem{Laxjoin} R.~Shibusawa, \emph{The join of two nuclei as a lax modality:
axiomatization, finite models, and inexpressibility}, preprint (2026),
archived at DOI 10.5281/zenodo.23330576.
\bibitem{Relnec} R.~Shibusawa, \emph{Joins of relativized necessities on a
poset: an alternating-chain criterion and the induction scheme of lax logic},
preprint (2026), archived at DOI 10.5281/zenodo.23411203.
\bibitem{Fuzzysame} R.~Shibusawa, \emph{Nuclei on the algebra of truth values
and the sameness lattice of fuzzy sets: meets, joins and the
finite-iteration axiom}, preprint (2026), archived at DOI 10.5281/zenodo.23413073.
\end{thebibliography}

\end{document}
